\section{Two-Level CBR: Level 1 --- Concrete Case Management}
\label{sec:level1}

\subsection{What a Level 1 Case Is}

A Level~1 case is not a data record about a past run. It is a \textbf{suspended
NormCode runtime} --- the full execution environment at a specific flow index, captured
and persisted.

\textbf{Definition.} A \textit{Level~1 case} $c$ is a tuple:
\begin{equation}
  c = (r,\; f,\; \mathcal{B}_f,\; \mathcal{C}_f,\; \mathcal{E}_f^{\,?})
\end{equation}
where $r$ is the run identifier; $f$ is the flow index of the last completed node;
$\mathcal{B}_f$ is the \textbf{Blackboard snapshot} at $f$: the full node status map
for all nodes at the moment $f$ completed; $\mathcal{C}_f$ is the \textbf{Concept
Repository snapshot} at $f$: all tensors for all nodes completed through $f$, indexed
by flow index; $\mathcal{E}_f^{\,?}$ is an \textbf{optional environment snapshot}:
present when $\mathcal{C}_f$ contains perceptual signs referencing external resources.

\textbf{Case quality guarantee.} Because the scope rule prevents any node from accessing
data outside its declared block, $\mathcal{C}_f$ is epistemically clean: every tensor
it contains is exactly the data that downstream nodes will use --- no hidden dependency
on accumulated context, no implicit reference to prior conversation history. Retrieving
$c$ will produce the same data-dependency behavior as the original execution,
\textit{provided} the external environment referenced by any perceptual signs in
$\mathcal{C}_f$ is accessible and consistent.

\subsection{The Case Base}

\textbf{Definition.} The \textit{case base} $\mathcal{CB}$ is the set of all Level~1
cases across all runs:
\begin{equation}
  \mathcal{CB} = \{(r, f, \mathcal{B}_f, \mathcal{C}_f, \mathcal{E}_f^{\,?}) \mid
  r \in \text{Runs},\; f \in \text{FlowIndices}(r)\}
\end{equation}

The case base is persisted in SQLite. The Checkpoint Panel in the Canvas App is the
visual interface to $\mathcal{CB}$.

\subsection{Retrieve}

Retrieval selects a past case from $\mathcal{CB}$ as the starting point for a new
execution. Two triggers: (1)~a run fails at flow index $f^*$; (2)~a new task is similar
to a past execution. Case inspection via the Tensor Inspector opens $\mathcal{C}_f$ at
any flow index in table, list, or JSON view. Because $\mathcal{C}_f$ contains exactly
the bounded inputs of each completed step, this inspection is O(1) --- C1.

\subsection{Reuse}

\textbf{Mechanism --- Fork.} The Fork operation implements reuse: (1)~new run identifier
$r'$ is created; (2)~orchestrator loads $(\mathcal{B}_f, \mathcal{C}_f)$ from retrieved
case; (3)~Blackboard marks all nodes at flow indices $\leq f$ as completed;
(4)~execution proceeds from the first pending node. No upstream re-execution.

\subsection{Revise}

\textbf{Mechanism --- Value Override + Selective Re-run:} (1)~operator retrieves a case
at flow index $f$; (2)~operator \textbf{overrides} a tensor in $\mathcal{C}_f$ via the
Tensor Inspector edit mode; (3)~orchestrator marks all downstream nodes as
\textbf{stale}; (4)~on resume, only stale nodes re-execute; (5)~revised run produces
new cases at every completed node.

\textbf{Scope-bounded revision (C3).} Because the scope rule guarantees that only nodes
downstream of $f$ can depend on the tensor at $f$, the stale set is exactly determined
by the data dependency graph. No manual specification required.

\subsection{Retain}

\textbf{Automatic checkpointing.} Every node completed during a run writes a new case
to $\mathcal{CB}$. Retention is fully automatic. The case base accumulates all runs'
Blackboard and Concept Repository snapshots continuously.

\subsection{Summary: The Level 1 CBR Cycle in Canvas}

\begin{table}[t]
\centering
\caption{Level~1 CBR operations mapped to Canvas mechanisms.}
\label{tab:level1cbr}
\begin{tabular}{@{}p{2.2cm}p{4.8cm}p{4.0cm}@{}}
\toprule
\textbf{CBR Operation} & \textbf{Canvas Mechanism} & \textbf{Key Property} \\
\midrule
Retrieve & Checkpoint Panel $\rightarrow$ select run/flow index; Tensor Inspector & O(1) inspection (C1) \\[4pt]
Reuse    & Fork: load $(\mathcal{B}_f, \mathcal{C}_f)$ into new run; no upstream re-execution & Isolated case = safe reuse \\[4pt]
Revise   & Value Override: edit tensor; selective re-run of stale boundary & Scope-bounded (C3) \\[4pt]
Retain   & Automatic checkpoint at every completed node $\rightarrow$ SQLite & Always, no operator action \\
\bottomrule
\end{tabular}
\end{table}
